\section{A receiver-centered framework for acoustic risk}
\label{sec:account}

The preceding sections do not converge on a single mechanism, but they do support a framework. Its premise is that acoustic meaning emerges from the interaction of signal information, prior experience and source reliability, ecological state and decision costs, and temporal history (\secref{sec:account-determinants}). Its organizing contribution is to recast those interactions as three recurring demands on the receiver---rapid acoustic defense, learned threat evaluation, and flexible action selection (\secref{sec:account-demands})---each of which draws on different evidence and makes different predictions. Because hearing, learned valuation, autonomic control, and defensive action run through partly separate pathways (\secref{sec:mechanisms}), the demands need not act on a single internal variable, and their effects can come apart. The framework is intended to organize evidence and generate tests, not to claim that a unified mechanism has been demonstrated.

\subsection{Four interacting determinants}
\label{sec:account-determinants}

How a receiver responds to a sound depends on four determinants (Sections~\ref{sec:features}--\ref{sec:persistence}): the \textbf{signal information} it can extract, its \textbf{prior experience} including the reliability of particular sources, its \textbf{ecological state and decision costs}, and the \textbf{temporal history} of exposure. What matters here is how they interact.

First, the determinants modify one another rather than adding up. Signal information is only as useful as experience makes it: rats freeze selectively to 22-kHz calls only after aversive experience, whereas naïve rats show slowing of movement and heart rate without that selective freezing \citep{KimEtAl2010,OlszynskiEtAl2020}. Experience is in turn source-specific, so a caller's history can make an otherwise effective alarm less effective in one species and more effective in another \citep{HareAtkins2001,BlumsteinEtAl2004}.

Second, the determinants act at different points. Ecological costs shape which action is expressed without necessarily changing whether a sound was recognized, as when cage familiarity switches mice between flight and freezing to the same sound \citep{MongeauEtAl2003}. Because hearing, valuation, autonomic control, and motor output run through partly separate pathways (\secref{sec:mechanisms}), a determinant can change one outcome---heart rate, say---while leaving another untouched.

Third, temporal history feeds back on all of them, through habituation, reliability learning, changing motivation, and possibly risk allocation, so that identical decline curves can conceal different processes (\secref{sec:persistence}). As in \figref{fig:framework}, these interactions form a receiver-centered model with feedback and a parallel route from detection to action.

\subsection{Three demands on the receiver}
\label{sec:account-demands}

The circuit literature suggests a sharper way to ask how the determinants combine. Consider three acoustic situations. A sudden loud or approaching sound can trigger flinching, freezing, or flight within a fraction of a second. A familiar conspecific alarm triggers defense because the receiver has learned what it predicts. And a repeated alarm from the same caller may lose or gain effect depending on who is calling, how reliable that caller has been, and what else is happening. These need not be three circuits. They are three demands---\emph{rapid acoustic defense}, \emph{learned threat evaluation}, and \emph{flexible action selection}---that place different requirements on the receiver and so make different predictions.

\textbf{Rapid acoustic defense} depends mainly on signal information. In laboratory mice and rats it can run through short pathways---cochlear root to PnC for startle, auditory cortex to IC and SC for flight---that do not require learned meaning \citep{DavisEtAl1982,LeeEtAl1996,XiongEtAl2015,LiEtAl2021}. Responses of this kind should be fast, present in inexperienced animals, and sensitive to level and onset; whether they also transfer across sources once acoustic features are matched is a hypothesis that has not been tested with natural calls. Fast is not the same as fixed: such responses still habituate, which is itself a form of learning \citep{RankinEtAl2009}, and state and context modulate them.

\textbf{Learned threat evaluation} depends mainly on experience and reliability. It requires the receiver to associate a sound with outcomes and to decide whether a new sound or caller belongs to a learned category. Laboratory candidates include thalamo-amygdala plasticity, the cortical--prelimbic interactions that limit generalization, and the hippocampal processes that maintain it \citep{LeDouxEtAl1984,QuirkEtAl1995,ConcinaEtAl2018,QuinnEtAl2009}. Experience-gated freezing to 22-kHz calls \citep{KimEtAl2010} and caller-specific discounting in ground squirrels \citep{HareAtkins2001} are behavioral instances. Responses of this kind should depend on history, be specific to sources or call types, and change when pairing with danger changes rather than with time alone.

\textbf{Flexible action selection} depends mainly on ecological state and costs. Once a sound has been evaluated, the receiver must still choose among vigilance, freezing, flight, and continued foraging. Competing CeA populations and parallel striatal and collicular outputs show how one evaluated threat can be routed to different actions \citep{FadokEtAl2017,TovoteEtAl2016,LiEtAl2021}; cage familiarity switching mice between flight and freezing to the same sound \citep{MongeauEtAl2003} and chipmunk responses that vary with refuge and food \citep{BaackSwitzer2000} are behavioral instances. Responses of this kind should change with refuge, food, and state even when the sound and the receiver's history are held constant, and autonomic signs of evaluation may persist when the overt action changes.

Temporal history acts on all three, but differently: rapid defense should habituate and recover with rest, learned evaluation should track reliability and source, and action selection should track changing costs. This is why decline curves alone cannot identify the process (\secref{sec:persistence}). The distinction is a way of comparing systems, not a claim that natural alarm calls engage any particular laboratory circuit. A single natural call might engage all three demands at once---an abrupt alarm from a familiar caller near a burrow, for example---which is precisely why experiments that vary one demand while holding the others fixed are informative.

These three demands are not exclusive to danger. A rat approaching a 50-kHz social call \citep{WilluhnEtAl2014} and a squirrel relaxing after hearing reassuring bird chatter \citep{LillyEtAl2019} must also recognize acoustic information, draw on experience, and select an appropriate response. The neural systems supporting attraction, safety, and defense need not be identical; the circuits reviewed in \secref{sec:mechanisms} are largely defensive, and the pup-call and dopamine findings show that other systems are involved when sounds signal opportunity \citep{MarlinEtAl2015,WilluhnEtAl2014}. What these examples share is the receiver's problem of deciding what a sound predicts and whether acting on it is worthwhile.

That problem can be written compactly as organizing notation. For the evaluated route, the receiver's choice resembles the cost-sensitive decision described in antipredator and escape theory \citep{LimaDill1990,YdenbergDill1986}:
\begin{equation}
a_t^{*}\in\arg\max_{a\in\mathcal{A}}\Big[\,p_t(D\mid x_t,k)\,\Delta L(a)\;-\;C(a\mid e_t,H_t)\Big].
\label{eq:action}
\end{equation}
Here $x_t$ is the acoustic information received at time $t$, and $k$ is the caller or source when the receiver can distinguish it. The term $p_t(D\mid x_t,k)$ is the receiver's current prediction that danger $D$ will follow, not a measurable ground truth. $\mathcal{A}$ is the set of feasible responses, from orienting and vigilance to freezing, escape, or continued feeding; $\Delta L(a)$ is the loss an action avoids relative to doing nothing; and $C(a\mid e_t,H_t)$ is its cost given current conditions $e_t$ (food, refuge, social setting, energetic state) and the recent history of danger and feeding $H_t$.

Each demand has its own place in this expression. Learned threat evaluation resides in $p_t$, which experience revises through prediction-error learning of the kind formalized by \citet{RescorlaWagner1972}, although field receivers often cannot observe whether a call was followed by a predator. Flexible action selection resides in $C$ and the choice among actions. Rapid acoustic defense sits outside the expression: fast, stimulus-locked responses interact or compete with the evaluated choice, as computational accounts of Pavlovian and goal-directed control describe \citep{DayanEtAl2006,RigoliEtAl2016}.

This expression is a heuristic, not a fitted model. No parameters are estimated, $a_t^{*}$ does not imply that animals compute an optimum, and similar decompositions already exist \citep{RigoliEtAl2016}; what is added here is their application to receivers of natural calls in laboratory and field. The notation makes one inference explicit: the same estimate of danger can produce different behavior when costs, available actions, or fast responses change, so similar behavioral curves need not share a mechanism. The accounts of decline in \secref{sec:persistence} act at different places---habituation on history-sensitive responses, including fast ones; reliability learning on $p_t$; motivational change and risk allocation on $C$---whereas temporal expectation needs more than a single probability to describe. Nor are these effects simple: unreliable marmot callers elicit more response, not less \citep{BlumsteinEtAl2004}, and habituation and risk allocation can interact in either direction \citep{Sirot2010}. The theories on which the notation draws are considered next.

\subsection{What existing theory explains---and where theories conflict}
\label{sec:account-theory}

Each body of theory captures part of the receiver's problem. Error management and signal detection explain why receivers accept many false alarms when misses are costly \citep{Nesse2005}, but not which features become predictive cues. Threat-sensitive defense explains graded responding to indicated risk \citep{Helfman1989} and fits urgency-linked vigilance, but it does not specify the mechanisms behind categorical escape or experience-gated freezing. The economics of escape explain why distance to refuge and lost foraging shape flight \citep{YdenbergDill1986,DillHoutman1989}, but most acoustic studies omit those variables. Predation risk allocation predicts less vigilance per alarm when high-risk periods are frequent \citep{LimaBednekoff1999}, a prediction that can look like habituation or unreliability learning in a decline curve.

Two further accounts reach toward mechanism. Predatory imminence organizes defensive modes by proximity of threat \citep{Fanselow1994}, connecting behavioral theory to the circuit families of \secref{sec:mechanisms-midbrain}. Computational accounts of Pavlovian and goal-directed control formalize how fast, stimulus-bound responses and evaluated choices compete \citep{DayanEtAl2006,RigoliEtAl2016}, but they were developed mainly from laboratory conditioning and human studies rather than natural calls. These theories complement one another; none says which computation operates in a given experiment, but the three demands above indicate where each applies.

\subsection{Principles emerging from the synthesis}
\label{sec:account-propositions}

Three principles emerge when the laboratory and field literatures are read together. None is a new mechanistic discovery; several of their components are well established within neuroscience or behavioral ecology. What is new is the question of whether these principles offer a common way to interpret natural acoustic communication across laboratory and field systems. Each principle makes predictions that could prove it wrong; \tabref{tab:predictions} collects the predictions that would discriminate between the processes involved. The questions they leave open are set out in \secref{sec:outlook}.

\subsubsection*{1. Acoustic structure constrains recognition, but experience determines its significance}

The useful question is not whether a natural alarm call outperforms a synthetic sound, but which acoustic features sensory processing privileges and how experience gives those features biological significance. Auditory physiology shows that feature coding comes before meaning: amygdala neurons in naïve rats separate 22-kHz from 50-kHz calls largely on frequency, without any freezing \citep{ParsanaEtAl2012a}, while abrupt or looming sounds reach defensive output through short brainstem and cortico-collicular routes that do not require learning \citep{DavisEtAl1982,XiongEtAl2015,LiEtAl2021}.

Experience then determines what those features come to mean. Naïve rats slow their movement and heart rate to 22-kHz calls, but selective freezing appears only after aversive experience \citep{OlszynskiEtAl2020,KimEtAl2010}, and predator calls give inconsistent results in predator-naïve laboratory stock \citep{HendrieNeill1991,KindermannEtAl2009}. In the field, gray squirrels respond to familiar heterospecific alarms and to unfamiliar alarms that resemble them, but less to unfamiliar dissimilar ones \citep{GetschowEtAl2013}. That is the pattern expected if learning operates on an acoustic similarity space that the auditory system supplies. Acoustic biases and learned associations interact; they are not competing ``innate'' and ``learned'' explanations.

This framing suggests why some sounds produce broadly generalized responses while others acquire meaning specific to a caller, context, or history. Features that act through rapid defensive routes---loudness, abrupt onset, looming, perhaps the harsh nonlinear components that heighten marmot responses \citep{BlumsteinRecapet2009}---should generalize widely, whereas features whose effects depend on learning should generalize along acoustic similarity only as far as circuits controlling generalization allow \citep{ConcinaEtAl2018,QuinnEtAl2009}. What the synthesis cannot yet explain is where any particular natural call falls between these poles. No study has asked whether the features that dominate auditory coding in a given species are also the features along which its learned alarm responses generalize. The principle predicts that responses carried by rapid routes will appear in naïve animals and transfer across acoustically matched callers, whereas learned responses will depend on history and generalize only along acoustic similarity; it would be challenged if a learned alarm response generalized as readily to acoustically dissimilar calls, or if a rapid-route response required prior experience (\tabref{tab:predictions}).

\subsubsection*{2. Recognizing danger and deciding what to do about it are separable problems}

The same evaluated threat can lead to different actions. Mice exposed to the same sweeps either freeze or flee depending on cage familiarity and prior shock \citep{MongeauEtAl2003}. Chipmunks hearing alarms become more vigilant, shorten their runs to cover, and delay re-emergence \citep{BaackSwitzer2000}, and gray squirrels hearing predator calls delay and reduce foraging without changing how far they stray from refuge \citep{JayneEtAl2015}.

The body can register a sound that behavior does not: naïve rats show cardiac slowing to 22-kHz calls without selective freezing \citep{OlszynskiEtAl2020}, juvenile ground squirrels discriminate alarm types by heart rate before they do so behaviorally \citep{Mateo1996}, and spiny mice raise cortisol to owl calls without visible defense \citep{EilamEtAl1999}. Laboratory circuits show how such separation can arise: cardiovascular and freezing components of a learned response travel through different CeA targets \citep{LeDouxEtAl1988}, competing CeA populations bias freezing versus flight \citep{FadokEtAl2017}, and parallel cortical outputs to striatum and colliculus produce freezing and flight to the same sound \citep{LiEtAl2021}. Escape theory, developed independently, predicts that the best action depends on refuge distance and lost foraging \citep{YdenbergDill1986}. Read together, these literatures suggest that behavior reflects an interaction between acoustic evaluation and the costs and opportunities of acting.

The consequence for interpretation is direct: a behavioral response cannot be read as a measurement of how threatening a sound is. A marmot that keeps feeding may have recognized the alarm and judged the costs of responding too high; a rat that freezes in an empty cage may have no better option. Comparisons of ``fear'' across species, settings, or call types are therefore confounded unless the options available to the animal are known. Whether sciurids use the circuits mapped in laboratory mice is unknown, and no field study has varied refuge or foraging opportunity while measuring an independent sign of recognition such as cardiac change. The principle predicts that such a manipulation will change the action a call elicits while leaving signs of recognition intact; it would be challenged if recognition measures shifted as much as behavior, which would imply that costs act on evaluation itself rather than on the choice of action (\tabref{tab:predictions}).

\subsubsection*{3. Experience updates different components of acoustic responding on different timescales}

A response to a repeated sound can decline for several reasons: the sound is no longer surprising, a particular caller has proven unreliable, the animal has learned when to expect it, or responding has become less worthwhile. Each process is documented somewhere. Source-specific decline with renewal after a caller change is established in ground squirrels and marmots \citep{Hare1998,BlumsteinDaniel2004}; discounting of callers whose alarms are not followed by danger is documented in juvenile ground squirrels \citep{HareAtkins2001}; irregular timing slows habituation of startle and of amygdala responses to artificial sounds \citep{Davis1970,HerryEtAl2007}; and approach to 50-kHz calls declines even when the calls predict social contact \citep{BerzEtAl2022}. Laboratory work on extinction and generalization offers candidate machinery for the source-specific components \citep{SierraMercadoEtAl2011,ConcinaEtAl2018}, and risk-allocation theory adds a slower, regime-level adjustment of vigilance to the overall frequency of danger \citep{LimaBednekoff1999}. The common thread is that experience can change how a sound is detected, what it predicts, or whether acting on it is worthwhile.

Habituation, associative learning, and risk allocation should therefore not be treated as competing labels for the same declining curve. In the notation introduced in \secref{sec:account-demands}, they act on different terms, and so predict different patterns of recovery, source specificity, and sensitivity to cost (\secref{sec:persistence-synthesis}). Caller-change recovery alone does not establish reliability learning, because it is also consistent with stimulus-specific habituation. Most studies measure only one of these signatures, so the relative contributions of each process to field responses remain unknown. The principle predicts that declines produced by different processes will part ways under rest, a change of caller, renewed pairing with danger, altered interval structure, and changed costs; it would be challenged if a single manipulation, such as rest alone, restored every declining response equally (\tabref{tab:predictions}).

% Main-text predictions table for Section 6.4; updating-process rows follow Figure 4 panels A--D.
\begin{table}[tbp]
\centering
\footnotesize
\singlespacing
\setlength{\tabcolsep}{5pt}
\renewcommand{\arraystretch}{1.15}
\setlength{\belowcaptionskip}{6pt}
\caption{Discriminating predictions of the receiver-centered framework. Rows group the three receiver demands (\secref{sec:account-demands}) and the four updating processes of \figref{fig:temporal}A--D. $p_t$ is the receiver's current prediction that danger will follow a sound; $C$ is the cost of an action given current conditions and recent history. No single prediction identifies a mechanism; convergent tests in the same animals are needed. The evidence column lists motivating studies and whether the prediction itself has been tested.}
\label{tab:predictions}
\begin{tabular}{@{}>{\raggedright\arraybackslash}p{1.05in}>{\raggedright\arraybackslash}p{0.9in}>{\raggedright\arraybackslash}p{2.05in}>{\raggedright\arraybackslash}p{2.05in}@{}}
\toprule
\textbf{Process} & \textbf{Where it acts} & \textbf{Discriminating prediction} & \textbf{Current evidence} \\
\midrule
\multicolumn{4}{@{}l}{\textit{Receiver demands}} \\
Rapid acoustic defense
  & Outside the expression (fast route)
  & Present in naïve animals and scales with level, onset, or looming. Hypothesis: transfers across callers when acoustic features are matched.
  & Synthetic sounds in laboratory rodents \citep{DavisEtAl1982,XiongEtAl2015,LiEtAl2021}; untested with natural alarms or acoustically controlled callers. \\
Learned threat evaluation
  & $p_t$
  & Depends on aversive or predator history; specific to call type or caller.
  & Experience-gated freezing to natural calls in rats; caller discounting in juvenile ground squirrels \citep{KimEtAl2010,HareAtkins2001}. \\
Flexible action selection
  & $C$ and the choice among actions
  & The action changes with refuge, food, or state while the sound, the receiver's history, and independent signs of recognition stay the same.
  & Context alters actions in mice and chipmunks \citep{MongeauEtAl2003,BaackSwitzer2000}; untested with an independent recognition measure. \\
\midrule
\multicolumn{4}{@{}l}{\textit{Updating processes}} \\
Habituation
  & History-sensitive responses, including fast ones
  & Partial recovery with rest alone; specific to the repeated stimulus, with transfer to other exemplars or callers in proportion to acoustic similarity.
  & Criteria defined \citep{RankinEtAl2009}; caller-specific decline and renewal in sciurids \citep{Hare1998,BlumsteinDaniel2004} fit habituation but do not identify the mechanism. \\
Reliability learning
  & $p_t$, by source
  & Decline specific to one source that reverses when its calls are again paired with danger, with acoustics and schedule held constant. The direction of change may differ between species.
  & Manipulated reliability changes responding in opposite directions in two sciurids \citep{HareAtkins2001,BlumsteinEtAl2004}; reversal untested. \\
Temporal expectation
  & Expected timing, beyond a single $p_t$
  & Slower decline with irregular than with regular intervals when mean interstimulus interval, stimulus properties, and number of exposures are matched.
  & Artificial sounds only \citep{Davis1970,HerryEtAl2007}; untested with natural calls. \\
Motivational and contextual relevance
  & $C$
  & Response magnitude or choice of action changes when the costs of responding change, with the sound held constant; recovery is one possible outcome.
  & Related observations in laboratory rodents \citep{MongeauEtAl2003,BerzEtAl2022}; untested with costs manipulated and recognition measured. \\
\bottomrule
\end{tabular}
\end{table}


\subsection{Limits of cross-species inference}
\label{sec:account-limits}

The framework connects two literatures; it does not by itself license moving a finding from one to the other. Three rules govern what can be carried across.

\textbf{Stimulus class.} Natural vocalizations, other natural sounds, and synthetic or conditioned sounds are different kinds of evidence (\secref{sec:mechanisms}). A circuit mapped with tones, noise, or looming sounds is a candidate mechanism for natural calls, not a demonstrated one, and a response to a simplified sound shows that its features suffice for that endpoint, not that natural calls carry nothing more.

\textbf{Evidence type.} Causal manipulations, activity measures, and studies of physiology or behavior alone support different claims. Activity during playback locates engagement without establishing necessity, and a physiological change without behavior shows that a sound was processed, not that it was evaluated as a threat.

\textbf{Species and setting.} Wild sciurids and laboratory rats and mice differ in repertoires, hearing range, ecological exposure, and behavioral options, and a laboratory cage offers neither food nor refuge. Field results establish what natural signals do under realistic costs; laboratory results establish mechanisms under controlled conditions. Neither substitutes for the other.

Applied together, these rules mark the central gap plainly. To our knowledge, no study has causally established the neural circuits mediating natural alarm-call responses in sciurids. Although playback studies have identified region-specific neural activation \citep{FreemanEtAl2020}, direct recordings or manipulations of defensive circuits during alarm-call responses in free-ranging sciurids remain unreported (\secref{sec:action-circuits}). The principles therefore say what to ask of a rat and a marmot side by side, and how cautiously to read the answers; they do not imply that both animals use the same circuits.

\subsection{Outstanding questions}
\label{sec:outlook}

Each principle above rests on field and laboratory evidence drawn from different animals, sounds, and measures, so none has yet been tested within a single system. Closing that gap defines three questions, one per principle, and a fourth about whether any of them holds beyond the few species and paradigms studied so far (the predictions that would answer the first three are in \tabref{tab:predictions}).

\textbf{Innate acoustic biases versus learning.} If acoustic structure constrains recognition while experience assigns significance, where does each natural call fall between broadly generalized and experience-specific responding? Predator-call studies in naïve laboratory rodents give conflicting answers that track strain, sound level, control stimulus, and endpoint \citep{HendrieNeill1991,KindermannEtAl2009}. Resolving the question requires following the same animals from first exposure through later experience, with autonomic and attentional measures alongside behavior, and asking whether unlearned biases shape what is later learned.

\textbf{Circuits for natural calls.} If recognition and action selection are separable, which pathways carry each when the sound is a natural call? Causal maps come largely from tones, noise, and synthetic looming sounds \citep{TovoteEtAl2016,FadokEtAl2017,XiongEtAl2015,LiEtAl2021}; for natural calls, the clearest causal evidence is the auditory thalamus requirement for experience-gated freezing to 22-kHz calls \citep{KimEtAl2010}. Answering the question requires pathway manipulations during natural-call playback, in settings where refuge, foraging, and caller history vary, and field measures of recognition that do not depend on overt defense. Animal-borne cardiac and neural loggers could supply those measures in free-living animals \citep{GaidicaDantzer2020,GaidicaDantzer2022}.

\textbf{Joint regulation by identity, reliability, and timing.} If experience updates different components of responding, how much of a given field decline does each process explain? Caller identity, reliability history, and timing have almost always been studied one at a time, and timing effects are known mainly from artificial sounds \citep{Hare1998,HareAtkins2001,Davis1970}. Resolving this requires designs that measure several signatures---recovery after rest, source specificity, reversal with renewed pairing, sensitivity to costs---in the same animals, with the number of calls, the mean interval, and the acoustic structure matched.

\textbf{Generalization across strains, populations, species, and contexts.} Most of the principles reviewed here rest on a few species and often a single paradigm each. Do they hold across laboratory strains, wild populations, and ecological settings? The evidence needed is comparative: the same acoustic manipulation and endpoints applied across species or populations under comparable costs, and laboratory effects re-examined once food and refuge are restored. The answer determines whether the synthesis describes shared mechanisms or only shared questions.
